\chapter{1977 Nobel Prize in Economics}
\textbf{Laureates: }Bertil Ohlin (Sweden), James Meade (UK)

\section*{Reason for Award}
For their pioneering contribution to the theory of international trade and international capital
movements

\section{What They Did}
Bertil Ohlin developed the Heckscher-Ohlin model, explaining trade patterns based on factor
endowments: countries export goods intensive in their abundant factors and import goods intensive
in their scarce factors. Ohlin's model, building on his teacher Eli Heckscher's earlier work,
provided a factor proportions theory of trade that explained why countries specialize according
to their relative factor abundances.

He also contributed to the theory of international capital movements, analyzing how capital flows
between countries in response to differential returns. Ohlin served as a Swedish political leader
and negotiator, applying his economic insights to international trade policy. He was involved in
the negotiation of the Treaty of Rome and served as a member of the Swedish parliament.

James Meade contributed fundamental theory of trade and customs unions, analyzing gains from
trade, tariff policies, and welfare effects of economic integration. Meade's analysis of customs
unions identified trade creation versus trade diversion effects, showing how preferential trade
agreements can both benefit and harm welfare depending on their structure.

He also contributed to the theory of international equilibrium, analyzing the balance of payments
and exchange rate determination. Meade's work on welfare economics provided the normative
foundation for evaluating trade policies. Together, their work provided a systematic framework
for understanding patterns of trade, factor price equalization implications, and the effects of
protectionism versus trade liberalization.

Meade's "The Theory of International Trade" (1957) provided a comprehensive treatment of
international trade theory and became a standard reference. His analysis of customs unions
demonstrated that preferential trade agreements could improve welfare through trade creation
while potentially reducing welfare through trade diversion. Ohlin's work on international trade
theory emphasized the role of factor endowments in determining trade patterns.

\section{Impact on Economic Thinking}
Their work provided a systematic framework for understanding patterns of trade and factor price
equalization implications. The Heckscher-Ohlin model explained why countries specialize and trade,
providing theoretical justification for reducing trade barriers. Meade's customs union analysis
clarified the welfare effects of economic integration in specific regions.

Their combined contributions shaped postwar trade liberalization efforts, influencing GATT and
WTO formation. The factor proportions theorem became central to trade theory, explaining pattern
differences between developed and developing nations. Their work also influenced debate about
protectionism, showing that while free trade is generally beneficial, specific policies might
improve welfare under certain conditions.

The Heckscher-Ohlin-Samuelson framework became standard in trade courses worldwide, teaching
students how relative factor abundance determines comparative advantage and trade flows. Meade's
welfare analysis provided the criteria for evaluating trade policies and international economic
integration. Their work established the theoretical foundation for analyzing the gains from trade
and the costs of protectionism that continues to inform trade policy debates.

Ohlin's political involvement demonstrated the practical application of economic theory to
policy-making. His work on international capital movements influenced understanding of financial
globalization and the balance of payments. Meade's analysis of customs unions provided the
theoretical foundation for understanding regional trade agreements and their welfare implications.

\section{Modern Perspectives Today}
The Heckscher-Ohlin framework has been extended incorporating differences in technology, human
capital, and transportation costs. New Trade Theory, incorporating economies of scale and
imperfect competition, complements the original model, informing analysis of global value chains,
multinational production, and services trade.

Patterns predicted by Heckscher-Ohlin explain broad trends in inter-industry trade, while New
Trade Theory explains intra-industry trade between similar countries. Together, these approaches
form the basis for modern trade policy negotiation and analysis. Empirical tests of factor
proportions predictions continue, with mixed results that have refined understanding of how trade
works in practice.

The model remains essential for analyzing trade agreements, tariff effects, and the distributional
consequences of globalization. Contemporary debates about offshoring, trade deficits, and
industrial policy build on foundations laid by Ohlin and Meade. Modern trade theory has
incorporated firm heterogeneity, showing that even within industries, only some firms export. The
analysis of global value chains and trade in tasks extends the traditional Heckscher-Ohlin
framework to contemporary realities of fragmented production.

Ohlin and Meade's contributions remain central to understanding the economic implications of
international trade and economic integration. Their work continues to inform trade policy debates
and the analysis of international economic relations. Modern empirical research continues to test
and refine the predictions of the Heckscher-Ohlin model, leading to a more nuanced understanding
of the factors determining trade patterns.
